%% kapa: exam=1 round=25 year=2014 period=1 form=A subject=civil_law subject_ko=민법 seq=36 number=36
%% source: Q-Net 감정평가사 기출문제 (cst00309) artlSeq=5101207 file=extracted/1008/5101207/2014년 제25회 감정평가사 1차시험 기출문제/1교시(A형).hwp
%% conversion: kordoc 4.24.51 -> tools/gen1.py (auto); review=proofread-sonnet; flags=none
\begin{question}{36}
민법상 유치권에 관한 설명으로 옳지 않은 것은? (다툼이 있으면 판례에 의함)
\begin{choices}
  \item 甲이 丙 소유의 자동차를 乙에게 수리시킨 경우, 乙은 그 수리비에 관하여 丙을 상대로 유치권을 주장할 수 있다.
  \item 채무자를 직접점유자로 하여 채권자가 간접점유하는 경우에도 유치권은 성립한다.
  \item 채권관계 당사자 사이의 그 채권에 관한 유치권 포기특약은 유효하다.
  \item 건물 신축공사의 하수급인이 다른 하수급인을 통하여 신축건물을 간접점유한 경우에도, 유치권의 성립요건을 충족할 수 있다.
  \item 공사대금채권에 기하여 유치권을 행사하는 자가 스스로 유치물인 주택에 거주ㆍ사용하는 것은 특별한 사정이 없는 한, 유치물의 보존에 필요한 사용에 해당한다.
\end{choices}
\end{question}
\answer{2}
